\section{Action weights in the three collision norms}
\label{sec:action-norm-details}
This section gives the norm calculation underlying
Proposition~\ref{prop:bounded-band-LY}. It also specifies the geometric
inputs when the scatterer is not a circle. The time parameter is fixed
throughout an orbit; a family of tables is not a randomly changing table.
We work on a local common space for an unforced finite-horizon periodic
billiard with strictly convex $C^3$ scatterers. Constants are uniform on
such a local family with common horizon, separation, curvature and
boundary bounds. A finite cover will suffice for a compact family.

\subsection{Norms, partitions and exponents}
We use densities relative to the invariant collision probability.
Thus a smooth function $h$ represents $h\nu$, and on densities
$\mathcal L h=h\circ T^{-1}$. This convention is important: the
Jacobian below is a stable-curve Jacobian, not an additional volume
Jacobian. Let $\mathcal W^s$ be the homogeneous stable curves of length
at most $\delta_0$, with the cone and curvature bounds in
\cite[Section 3.3]{DZ}. Locally write $W=\{(r,\varphi_W(r)):r\in J_W\}$
and let $G_W$ be this graph parametrization. Curves in different
homogeneity strips have infinite distance; within one strip put
\[
 d_s(W_1,W_2)=|J_{W_1}\mathbin\triangle J_{W_2}|
       +\|\varphi_{W_1}-\varphi_{W_2}\|_{C^1(J_{W_1}\cap J_{W_2})}.
\]
The spaces $C^a(W)$ below are the closures of smooth functions in the
$a$-H\"older norm. For smooth $h$ set
\begin{align}
 |h|_w&=\sup_{W\in\mathcal W^s}\sup_{\|\psi\|_{C^p(W)}\le1}
                         \left|\int_W h\psi\,dm_W\right|,
                                             \label{eq:v35-weak-norm}\\
 \|h\|_s&=\sup_{W\in\mathcal W^s}
       \sup_{\|\psi\|_{C^q(W)}\le |W|^{-\varsigma}}
                         \left|\int_W h\psi\,dm_W\right|,
                                             \label{eq:v35-stable-norm}\\
 \|h\|_u&=\sup_{0<\epsilon\le\epsilon_0}
       \sup_{d_s(W_1,W_2)\le\epsilon}
       \sup_{\substack{\|\psi_i\|_{C^p(W_i)}\le1\\
                \psi_1\circ G_{W_1}=\psi_2\circ G_{W_2}}}
       \epsilon^{-\gamma}
       \left|\int_{W_1}h\psi_1\,dm_{W_1}
                    -\int_{W_2}h\psi_2\,dm_{W_2}\right|.
                                             \label{eq:v35-unstable-norm}
\end{align}
Equality of the graph tests is imposed on the common interval. The
strong and weak spaces are the completions of $C^1$ densities in
$\|h\|_s+c_0\|h\|_u$ and $|h|_w$, respectively. Changing $c_0>0$
changes the norm, not its completion. These are the conventions in
\cite[Section 3.1]{DPZ}.

The following strict choices are compatible with all the estimates used
here:
\begin{equation}\label{eq:v35-exponent-choice}
 p=\frac1{12},\qquad q=\frac1{24},\qquad
 \varsigma=\frac18,\qquad \zeta=\frac14,\qquad
 0<\gamma\le\frac1{100}.
\end{equation}
If necessary decrease $\gamma$ to meet the local geometric restriction.
For the usual homogeneity strips the inverse-Jacobian summability
threshold can be chosen as $\zeta_0=2/3$. In particular
\[
 \varsigma<1-\zeta_0,\quad
 \gamma<\min\{\varsigma,p-q,1/3-q,1-q\},\quad
 \max\{p,\gamma/(1-\gamma)\}<\zeta<1/2.
\]
The threshold statement uses only the standard quadratic homogeneity
strips and the one-step summability in \cite[Section 3.1]{DZ}.
The first two restrictions are the norm restrictions; the last is the
piecewise multiplier restriction. No equality at a H\"older endpoint is
used.

Here is the precise partition condition on a preceding-flight record
$\check f=f\circ T^{-1}$. In each homogeneity strip its regularity
partition has at most $N_0$ simply connected pieces. A stable curve
intersects any piece in at most $K_0$ intervals, and
\begin{equation}\label{eq:v35-partition-boundary}
 m_W\{x\in W:d(x,\partial Z)<\epsilon\}\le C_0\epsilon
\end{equation}
for each partition piece $Z$. The constants are common to all strips
and local parameters. On each piece $\check f$ has a common
$C^{1/2}$ bound. A finite refinement by coordinate cuts is permitted,
but a number of cuts growing with an orbit length is not part of this
condition. The circular root calculation in
Lemma~\ref{lem:backward-flight-multiplier} verifies these conditions;
Lemma~\ref{lem:v35-ellipse-geometry} will verify them for ellipses.

\begin{lemma}[The precise multiplier consequence]
\label{lem:v35-multiplier-detail}
Under this partition condition, $M_{\exp(iz\cdot\check f)}$ is an
entire strong-space multiplier. For a compact complex set $K$ and any
multi-index $\beta$,
\begin{equation}\label{eq:v35-multiplier-derivatives}
 \sup_{z\in K}\|\partial_z^\beta
        M_{\exp(iz\cdot\check f)}\|_{\mathcal B\to\mathcal B}
       \le C_{K,\beta}.
\end{equation}
The same bound holds on the weak space.
\end{lemma}
\begin{proof}
The conditions just listed are precisely the bounded-complexity
conditions of \cite[Lemma 3.3]{DPZ}; its multiplier conclusion applies
because of \eqref{eq:v35-exponent-choice}. Products on each piece obey
$\|uv\|_{C^\zeta}\le\|u\|_{C^\zeta}\|v\|_{C^\zeta}$.
If the components of $\check f$ have norm at most $F$, the exponential
series in that algebra is bounded by $\exp(CF|z|)$. Its derivative is
multiplication by $(i\check f)^\beta\exp(iz\cdot\check f)$ and has
norm at most $(CF)^{|\beta|}\exp(CF|z|)$. The fixed partition has not
changed. Applying the multiplier bound to the series and its locally
uniformly convergent derivatives proves the strong assertion and
analyticity. For the weak norm one splits the stable test at the same
finite boundaries and uses its $C^p$ bound; alternatively the weak part
of the same multiplier estimate gives the assertion. The local boundary
reparametrizations have uniformly bounded derivatives and inverses,
so the partition and H\"older bounds survive them.
\end{proof}

\subsection{The geometric sums, before inserting a weight}
Write $\mathcal G_m(W)$ for the homogeneous inverse pieces, subdivided
at the reference length as in \cite[Section 3.2]{DZ}, and let
$\mathcal I_m(W)$ be the pieces without a long ancestor. For
$U\in\mathcal G_m(W)$ set
$j_U=\|J_UT^m\|_\infty$. The following are the unweighted geometric
inputs, recorded with their length factors:
\begin{align}
 \sum_{U\in\mathcal G_m(W)}j_U&\le C_2,&
 \sum_{U\in\mathcal G_m(W)}
                 (|U|/|W|)^a j_U&\le C_2^{1-a}
                   &&(0\le a\le1),                 \label{eq:v35-geometric-sums-a}\\
 \sum_{U\in\mathcal I_m(W)}j_U&\le C_1\theta_*^m,&
 \sum_{U\in\mathcal G_m(W)}j_U^{1-\varsigma}&\le C_3^m,
                                                        \label{eq:v35-geometric-sums-b}
\end{align}
where $\theta_*<1$. Precisely, the never-long sum is \cite[Lemma 3.2(a)]{DZ}, the total
sum is its part (b), the length-weighted sum is part (c), and the
inverse-Jacobian power sum is part (d), on page 12 of the pinned
arXiv:1210.1261v1 text. The
last inequality uses $1-\varsigma>\zeta_0$. Bounded distortion gives
$\|J_UT^m/j_U\|_{C^p(U)}\le C$. Jensen's inequality and
$\sum_U j_U|U|\le C|W|$ also give
\begin{equation}\label{eq:v35-never-long-length-sum}
 \sum_{U\in\mathcal I_m(W)}(|U|/|W|)^{\varsigma}j_U
                         \le C\theta_*^{(1-\varsigma)m}.
\end{equation}
We retain these sums as geometric theorems, not consequences of a
finite-word computation. All their constants are independent of a
Fourier frequency.

The symplectic input needed to insert the roof is
\begin{equation}\label{eq:v35-action-general}
 d\tau=T^*\vartheta-\vartheta,\qquad \vartheta=\sin\varphi\,dr.
\end{equation}
It follows by differentiating the length of a regular free flight at
its two endpoints. The tangential component of the velocity is unchanged
by reflection. Thus the derivative at the terminal endpoint is the
outgoing tangential momentum there, and the derivative at the initial
endpoint is its negative. This proves the formula also for noncircular
scatterers.

For real $(u,b)$ with $|b|\le B$ put
$w_m=\exp\{i(u\cdot K_m+b(S_m-m\bar\tau))\}$.
On a regular branch $K_m$ is constant. For an arc $\gamma$ joining
$x$ to $y$ and regular through time $m$, integration of
\eqref{eq:v35-action-general} gives
\begin{equation}\label{eq:v35-action-telescope}
 S_m(y)-S_m(x)=\int_{T^m\gamma}\vartheta-\int_\gamma\vartheta.
\end{equation}
Uniform boundedness of $\vartheta$ and stable contraction therefore
bound $\|w_m|_U\|_{C^1}$ by $C_B$, independently of $m$.

For completeness the matching conditions used here are the following.
The untrimmed inverse graphs are joined by the unstable-cone foliation
of \cite[Lemma 3.3 and Section 4.3]{DZ}. Every intermediate connector
and both graphs lie in one regular branch and one homogeneity strip;
a connector cut by a singularity is placed in the unmatched family.
The construction gives
\[
 |\gamma_x|\le C\Lambda^{-m}\epsilon,\quad
 |T^m\gamma_x|\le C\epsilon,\quad
 |T^i\gamma_x|\le2\max_{j=1,2}|T^iU_j|\quad(0\le i\le m).
\]
In particular the two endpoints have the same lattice itinerary. These
are the hypotheses of the matched Jacobian estimate, rather than an
additional claim about all nearby orbits. Trimming two such graphs to
the common base interval moves an endpoint by at most $C\epsilon$
on a stable graph, and its image by at most $C\Lambda^{-m}\epsilon$.
Equation~\eqref{eq:v35-action-telescope} consequently gives
\begin{equation}\label{eq:v35-weight-match-detail}
 \|w_m^{(1)}-w_m^{(2)}\|_\infty\le C_B\epsilon,
 \qquad
 \|w_m^{(1)}-w_m^{(2)}\|_{C^q}\le C_B\epsilon^{1-q}.
\end{equation}
The latter follows by interpolation with the individual Lipschitz
bounds in common graph coordinates. It is valid at every length $m$.
No forward flight singularity or homogeneity boundary has been removed:
all such cuts are already in the inverse-curve partition. The weight
adds none because its only discrete label is constant on a regular
collision branch.

\subsection{The weighted inequalities with all length factors}
\begin{proposition}[Action-weighted collision inequalities]
\label{prop:v35-expanded-LY}
For each fixed $B<\infty$, the exact operator
$Q_\omega=M_{\exp(i\omega\cdot\check f^{\rm c})}\mathcal L$ satisfies
\begin{align}
 |Q_\omega^m h|_w&\le C_B|h|_w,                  \label{eq:v35-LY-w}\\
 \|Q_\omega^m h\|_s
 &\le C_B\bigl(\Lambda^{-qm}
             +\theta_*^{(1-\varsigma)m}\bigr)\|h\|_s+C_B|h|_w,
                                                       \label{eq:v35-LY-s}\\
 \|Q_\omega^m h\|_u
 &\le C_B\Lambda^{-\gamma m}\|h\|_u+C_BC_3^m\|h\|_s.
                                                       \label{eq:v35-LY-u}
\end{align}
Here $C_3$ may be enlarged. The exponents and geometric constants are
independent of $B$; the multiplier constants are not.
\end{proposition}
\begin{proof}
It suffices first to use smooth densities. A pairing on a terminal
curve is the sum of
\[
 \int_U h\,w_m J_UT^m(\psi\circ T^m)\,dm_U.
\]
The normalized Jacobian and $w_m$ have uniformly bounded $C^p$ norms,
where the latter bound is $C_B$. The first sum in
\eqref{eq:v35-geometric-sums-a} proves \eqref{eq:v35-LY-w}.

For the stable norm let $a_U$ be the arclength average of
$\psi\circ T^m$ on $U$. Then
\[
 \|\psi\circ T^m-a_U\|_{C^q(U)}
             \le C\Lambda^{-qm}|W|^{-\varsigma}.
\]
The contribution of these differences is bounded by
\[
 C_B\Lambda^{-qm}\|h\|_s
       \sum_U (|U|/|W|)^{\varsigma}j_U
       \le C_B\Lambda^{-qm}\|h\|_s.
\]
On never-long pieces the average contribution is bounded by
$C_B\|h\|_s$ times \eqref{eq:v35-never-long-length-sum}.
For the remaining pieces group by the most recent long ancestor
$V\in\mathcal G_k(W)$, with $|V|\ge\delta_0/3$. Descendants with no
later long ancestor have Jacobian sum at most $C\theta_*^{m-k}$
relative to $V$. The smooth average, weight and normalized Jacobian
are paid by the weak norm, exactly as for the average term in
\cite[equations (4.9)--(4.11)]{DZ}. Keeping the length normalization,
the resulting bound is at most
\[
 C_B\delta_0^{-\varsigma}|h|_w
   \sum_{k=0}^m\theta_*^{m-k}
     \sum_{V\in\mathcal G_k(W)} (|V|/|W|)^{\varsigma}j_V
 \le C_B|h|_w.
\]
The $k=0$ term is included only if the root is long. This proves
\eqref{eq:v35-LY-s}. In particular the noncontracting derivative of
$w_m$ has not been multiplied by $\|h\|_s$ in its leading term.

We spell out which unstable errors the last geometric sum pays.
An unmatched piece $V$ has $|T^mV|\le C\epsilon$; bounded distortion
implies $j_V|V|\le C\epsilon$. Its stable-norm contribution is at most
$C_B j_V|V|^{\varsigma}\|h\|_s$. Hence all unmatched pieces cost at most
\[
 C_B\epsilon^{\varsigma}\|h\|_s
                 \sum_V j_V^{1-\varsigma}
       \le C_B\epsilon^{\varsigma}C_3^m\|h\|_s.
\]
After division by $\epsilon^\gamma$ this is admissible.
On a matched pair the old graph/test comparison, with one weight
attached, has the leading term
$C_B\epsilon^\gamma\Lambda^{-\gamma m}\|h\|_u$.
The normalized Jacobian and graph reparametrization errors have
$C^q$ size bounded by
$C(\epsilon^{p-q}+\epsilon^{1/3-q})$ in the matching calculation of
\cite[Section 4.3]{DZ}. Their stable-norm sums, with each Jacobian
factored out, are bounded by $C C_3^m\|h\|_s$.
The additional difference of the two weights has $C^q$ size
$C_B\epsilon^{1-q}$ by \eqref{eq:v35-weight-match-detail}, and the
same Jacobian/length sum bounds its stable contribution. Each of
$\varsigma,p-q,1/3-q,1-q$ exceeds $\gamma$.
Dividing the total difference by $\epsilon^\gamma$ proves
\eqref{eq:v35-LY-u}. This calculation retains the exponential stable
cross coefficient. Completion of smooth densities gives the result.
\end{proof}

Fix $B$ before making the following choices. Choose $N$ so that both
leading coefficients in \eqref{eq:v35-LY-s} and
\eqref{eq:v35-LY-u} are less than $1/4$. Then choose $c_B>0$ so that
$c_B C_B C_3^N<1/4$. In the equivalent norm
$\|h\|_B=\|h\|_s+c_B\|h\|_u$,
\[
 \|Q_\omega^N h\|_B\le\tfrac12\|h\|_B+C'_B|h|_w.
\]
Iteration and \eqref{eq:v35-LY-w}, followed by the finitely many
remaining powers, prove power boundedness. Compact embedding of the
strong unit ball in the weak space gives a common essential radius
less than one for this fixed band. No estimate for the growth of
$N,c_B^{-1}$ or the spectral constants as $B\to\infty$ is implied.

\subsection{From a peripheral distribution to a physical phase}
We isolate a small functional-analytic detail. The weak completion
alone must not be identified with a faithful space of global
functionals without checking that identification.

\begin{lemma}[Faithfulness on the strong completion]
\label{lem:v35-faithful-strong}
If $h\in\mathcal B$ vanishes on all smooth global tests, then $h=0$.
\end{lemma}
\begin{proof}
Fix a stable graph compactly inside one homogeneity strip and a smooth
test supported in its interior. Its small transverse translates are
admissible stable graphs after choosing the cone and curvature bounds
with a strict margin. Give them the same graph-coordinate test. The
unstable norm bounds the difference of their pairings by
$C\|h\|_u|v|^\gamma$, where $v$ is the transverse translation.
Average these pairings against a smooth transverse kernel of width
$\epsilon$. On smooth densities the average is a global test pairing:
the graph arclength Jacobian is included in the test and the collision
density is divided out. This density is smooth and bounded below on
the compact strip under consideration. The resulting test is $C^1$;
smooth approximation and the weak pairing bound permit its use for
$h\in\mathcal B$. By the hypothesis the average is zero. Letting
$\epsilon\downarrow0$ makes the original graph pairing zero.

General little-H\"older tests are limits of smooth tests. Trimming
endpoints has error at most $C\|h\|_s\ell^{\varsigma-q}$ on a trimmed
interval of length $\ell$. Curves on a boundary of the admissible
family can first be approximated by ones with strict cone and curvature
margins, using the unstable comparison and the same trimming estimate.
We now conclude in the strong norm, without using an injective inclusion.
For a fixed admissible $W$, its pairing is continuous in the little
$C^q(W)$ norm by \eqref{eq:v35-stable-norm}. Approximate each such test
in that norm by smooth tests, after the endpoint and curve
approximations just described. The pairings therefore vanish for every
$C^q$ stable test, not merely for the $C^p$ tests of the weak norm.
Their normalized supremum is zero, so $\|h\|_s=0$.
Every difference in \eqref{eq:v35-unstable-norm} is now a difference
of two zero pairings. For each fixed pair and each fixed $\epsilon>0$
its normalized value is zero. Taking the defining supremum gives
$\|h\|_u=0$.

These equalities also hold for the completed, rather than just the
smooth, norm. Indeed the defining normalized stable pairings and
normalized pair differences embed smooth densities isometrically in
the product of two spaces of bounded scalar families. A strong Cauchy
sequence converges uniformly in those two families. Thus their
suprema give the completed seminorms as well; no interchange of an
uncontrolled pointwise limit with a supremum is involved. It follows
that $\|h\|_s+c_0\|h\|_u=0$ and hence $h=0$ in the completion.
This argument uses transverse regularity of the strong space and
asserts no faithfulness for arbitrary weak-completion vectors.
The independent inclusions recorded in \cite[Lemma 3.2(c), footnote 7]{DPZ}
are consistent with, but are not used to prove, this conclusion.
\end{proof}

\begin{lemma}[Bounded physical representatives of peripheral vectors]
\label{lem:v35-peripheral-density}
Let $Q_\omega$ be power bounded and quasi-compact as above. If
$|\lambda|=1$ is an eigenvalue, a nonzero peripheral vector is represented
by $q\nu$ with $q\in L^\infty(\nu)$. After multiplication by a constant,
$|q|=1$ and
\begin{equation}\label{eq:v35-peripheral-physical}
 q\circ T=\lambda^{-1}e^{i\omega\cdot f^{\rm c}}q
                    \quad\hbox{almost everywhere}.
\end{equation}
\end{lemma}
\begin{proof}
Power boundedness excludes Jordan blocks at modulus one. The Cesaro
averages $N^{-1}\sum_{j<N}\lambda^{-j}Q_\omega^j$ converge in the strong
space to the spectral projection $P_\lambda$. Since $C^1$ densities
are dense and its range is finite dimensional, $P_\lambda g\ne0$ for
some smooth bounded density $g$ (not necessarily positive). On physical
densities every summand has modulus at most $\|g\|_\infty$. Thus for
smooth global $\phi$,
\[
 |(P_\lambda g)(\phi)|\le\|g\|_\infty\|\phi\|_{L^1(\nu)}.
\]
Extension by density in $L^1$ supplies $q\in L^\infty(\nu)$. It is
nonzero by Lemma~\ref{lem:v35-faithful-strong}.

To justify the phase equation without testing an abstract distribution
on a discontinuous pullback, regard the same averages as bounded
physical densities. Their weak-star limits are precisely $q$, because
smooth global tests are dense in $L^1$. Weighted composition by $T^{-1}$
is weak-star continuous on $L^\infty$, with its usual $L^1$ predual.
The difference between its action on the average and $\lambda$ times
the average has supremum norm at most $2\|g\|_\infty/N$. Taking the
weak-star limit proves the physical eigen-equation and hence
\eqref{eq:v35-peripheral-physical}. Ergodicity makes $|q|$ a positive
constant. Division by that constant finishes the proof.
\end{proof}

The arithmetic result in the next section now applies to an actual
measurable function, not an anisotropic distribution. Exclusion of
nonzero peripheral frequencies, local strong-to-weak continuity,
and the fixed-band inequalities yield a uniform spectral gap on each
compact set avoiding zero. For parameter continuity the one-step
comparison uses the same matched and unmatched curves as
\cite[Section 5]{DZ}. On matched pieces the preceding label agrees and
the explicit root is jointly $1/2$-H\"older in state and parameter;
interpolation gives convergence in $C^q$. The unmatched short pieces
are bounded by the same small-length sum. This explains exactly where
the parameter-uniform preceding-root condition is used. No continuity
of a long singular itinerary is required.
